\chapter{LITERATURE SURVEY}

\section{Classical Image Enhancement Approaches}
Classical low-light enhancement techniques relied primarily on spatial and histogram manipulation:
\begin{itemize}
    \item \textbf{Histogram Equalization (HE) and CLAHE}:
    Contrast-Limited Adaptive Histogram Equalization (CLAHE) divides the image into contextual tiles and equalizes local contrast while limiting noise amplification via slope clipping. While computationally lightweight ($< 5\text{ ms}$), it operates blindly on pixel intensities: in extreme low light, it severely amplifies background shot noise and causes harsh edge halos.
    \item \textbf{Retinex Theory (Land, 1977)}:
    Retinex theory \cite{land1977retinex} decomposes an observed image $S$ into illumination $L$ and reflectance $R$ via:
    \begin{equation}
    S(x, y) = L(x, y) \cdot R(x, y)
    \label{eq:retinex}
    \end{equation}
    Classical variational Retinex methods (Single-Scale Retinex, Multi-Scale Retinex) assume illumination varies smoothly across space while reflectance contains textural details. These methods require solving expensive optimization problems per image and consistently wash out colors in dark regions.
\end{itemize}

\section{Deep Learning in the sRGB Domain}
With the advent of deep convolutional neural networks, researchers trained models to estimate Retinex components or contrast curves:
\begin{itemize}
    \item \textbf{Retinex-Net (Wei et al., 2018)}:
    Combines a Decom-Net (to separate illumination and reflectance) with an Enhance-Net and BM3D denoiser \cite{wei2018deep}. Because it operates on compressed 8-bit sRGB images, severe color distortion and blotchy noise artifacts appear whenever dark areas are amplified.
    \item \textbf{Zero-DCE (Guo et al., 2020)}:
    Formulates low-light enhancement as a task of estimating pixel-wise, high-order curve parameter maps \cite{guo2020zero}. It requires no paired training data. While extremely fast ($0.08\text{M}$ parameters, $> 100\text{ FPS}$), Zero-DCE has no denoising mechanism; boosting curves in extreme darkness blows up shot noise into an unreadable mess ($18.20\text{ dB}$ PSNR).
\end{itemize}

\section{Raw Domain Deep Learning}
\begin{itemize}
    \item \textbf{Learning to See in the Dark (Chen et al., 2018)}:
    Chen et al. \cite{chen2018learning} made the seminal breakthrough of training a fully convolutional U-Net directly on packed Bayer raw data ($H/2, W/2, 4$) from a Sony $\alpha$7S II sensor. Bypassing demosaicing allowed the network to learn joint denoising, demosaicing, and color correction from paired short/long exposures. They proved that raw data contains vast amounts of recoverable detail that traditional ISPs discard.
    
    \textit{Critical Shortcomings}:
    \begin{enumerate}
        \item The model is static. It accepts an amplification factor $\alpha$ as a fixed scalar input multiplied directly into the input pixels ($X_{in} = \alpha \cdot X$). If the user alters $\alpha$, the network frequently hallucinates strange color shifts because the convolutional weights were trained on a fixed ratio.
        \item Standard skip connections allow high-frequency noise from shallow layers to corrupt the decoder output.
        \item It requires $> 120\text{ ms}$ per frame on desktop GPUs and $> 140\text{ ms}$ on mobile CPUs, rendering it non-viable for real-time mobile viewfinders.
    \end{enumerate}
\end{itemize}

\section{Generative Models and Diffusion LLIE Paradigms}
\begin{itemize}
    \item \textbf{PatchGAN and Pix2Pix (Isola et al., 2017)}:
    Isola et al. \cite{isola2017image} established that standard $L_1$ and $L_2$ losses produce blurry results because they average all plausible high-frequency solutions. Combining $L_1$ with an adversarial PatchGAN discriminator penalizes structure at the scale of local image patches, forcing the generator to reconstruct crisp textures.
    \item \textbf{Diffusion Models (e.g., DiffLight, 2024)}:
    Denoising Diffusion Probabilistic Models (DDPM) model image distribution through iterative Langevin dynamics. While DiffLight achieves high PSNR ($29.10\text{ dB}$) and realistic fine textures, the reverse diffusion process requires $15\text{ to }50$ sampling steps. Even on high-end desktop GPUs, latency exceeds $3.4\text{ seconds}$ per frame. On mobile hardware, diffusion models overheat the device and consume excessive battery power, rendering them non-viable for real-time camera viewfinders.
\end{itemize}

\section{Edge Quantization and Mobile Acceleration}
\begin{itemize}
    \item \textbf{Post-Training Quantization (PTQ) (Jacob et al., 2018)}:
    Deploying multi-million parameter neural networks onto mobile devices requires model compression \cite{jacob2018quantization}. PTQ maps 32-bit floating-point weights and activations ($\mathbb{R}$) to 8-bit signed integers ($\mathbb{Z} \in [-128, 127]$) via uniform affine mapping:
    \begin{equation}
    q = \text{round}\left(\frac{x}{S}\right) + Z
    \label{eq:ptq_mapping}
    \end{equation}
    where $S$ is the scale factor and $Z$ is the zero-point. INT8 operations execute on mobile Neural Processing Units (NPUs) and Digital Signal Processors (DSPs) using integer SIMD dot-product instructions, reducing memory bandwidth by $4\times$ and latency by $2\text{ to }4\times$ compared to FP32 CPU execution.
\end{itemize}

\section{Comprehensive Gap Analysis Matrix}
Table \ref{tab:gap_analysis} summarizes the comparative analysis of existing paradigms versus the proposed LumaCam architecture.

\begin{table}[t]
\centering
\caption{Comparative Analysis: Existing LLIE Paradigms vs. Proposed LumaCam Architecture}
\label{tab:gap_analysis}
\resizebox{\textwidth}{!}{%
\begin{tabular}{lcccccc}
\toprule
\textbf{Method / Paradigm} & \textbf{Operating Domain} & \textbf{PSNR (dB)} & \textbf{SSIM} & \textbf{Latency (ms)} & \textbf{Interactive $\alpha$?} & \textbf{Mobile Edge Viable?} \\ \midrule
Linear Digital Gain ($\times 100$) & Bayer RAW & 14.12 & 0.385 & $< 1.0$ & Static Linear & No (Severe Noise Blowup) \\
Retinex-Net (Wei et al., 2018) \cite{wei2018deep} & sRGB JPEG & 22.40 & 0.680 & $\approx 65.0$ & No & Marginal (Color Distortions) \\
Zero-DCE (Guo et al., 2020) \cite{guo2020zero} & sRGB JPEG & 18.20 & 0.560 & $\mathbf{8.5}$ & Curve Param & Poor (No Denoising Ability) \\
Chen et al. SID (CVPR 2018) \cite{chen2018learning} & Bayer RAW & 28.62 & 0.792 & $142.0$ & Static Fixed & No ($> 140\text{ms}$ on Mobile CPU) \\
DiffLight (Diffusion, 2024) & Bayer RAW & $\mathbf{29.10}$ & 0.805 & $3,450.0$ & Iterative Steps & No ($3.45\text{s}$ Desktop Latency) \\ \midrule
\textbf{LumaCam (FP32 Core)} & \textbf{Bayer RAW} & \textbf{28.88} & $\mathbf{0.812}$ & $\mathbf{18.6}$ (DirectML) & \textbf{Yes (FiLM Latent)} & \textbf{Yes (35--50 FPS)} \\
\textbf{LumaCam (INT8 Quantized)} & \textbf{Bayer RAW} & \textbf{28.56} & \textbf{0.806} & $\mathbf{24.8}$ (Pixel 7 NPU) & \textbf{Yes (FiLM Latent)} & \textbf{Yes (Edge-Native 40.3 FPS)} \\ \bottomrule
\end{tabular}%
}
\end{table}

\section{Inferences and Research Challenges}
From the literature review and gap analysis, three principal research challenges emerge:
\begin{enumerate}
    \item \textbf{The Demosaicing Paradox}: Performing denoising after demosaicing is mathematically suboptimal because noise correlation is already scrambled across channels. Restoration must occur in the raw Bayer CFA domain.
    \item \textbf{The Rigidity Challenge}: Current deep raw restoration networks lack real-time exposure adaptability. Users require a tactile slider ($\alpha$) that modulates deep latent representations without triggering color distortions.
    \item \textbf{The Edge Constraint}: State-of-the-art diffusion models cannot meet the real-time mobile threshold ($\le 33.3\text{ ms}$). A single-pass Attention U-Net GAN, optimized via INT8 Post-Training Quantization, is essential to bridge the efficiency-quality gap.
\end{enumerate}
